\chapter{Open Questions Register}
\label{app:open-questions}

\begin{quotation}
\noindent\textit{Synopsis.} This appendix maintains a single consolidated
record of every conjecture and open question raised in the body, keeping a
clear separation between established results and speculative extensions.
\end{quotation}

Per the epistemic-labeling discipline stated in the front matter, this
appendix consolidates every claim in the body labeled Conjecture or Open
Question, together with every substantive open item flagged informally in
a Remark rather than housed in a formal environment, kept separate from
proved content so that the monograph's unproved commitments are never
left implicit. Entries are grouped by the part of the framework they
bear on, rather than by chapter order, since several recur across
chapters far apart in the text.

\begin{itemize}
  \item \textbf{The MEM|8 fork} (\cref{oq:mem8-fork}, Ch.~\ref{ch:mem8}).
  The single largest open dependency of the monograph: whether MEM|8's
  claim generalizes from biological recall specifically (narrow reading)
  to representation as such (broad reading), and, if the latter, whether
  Part~\ref{part:foundations} requires revision from the ground up.
  \cref{oq:unmediated-access} (Ch.~\ref{ch:fidelity-continuum}) is the
  same fork applied to the specific question of whether any instance of
  ``direct access'' involves strictly zero mediation.

  \item \textbf{Discard-rule robustness --- partially resolved}
  (Ch.~\ref{ch:budget-relay}, \cref{rem:discard-rules-diverge};
  Appendix~\ref{app:proof}). \cref{def:budget-discard} is satisfied by
  more than one non-equivalent construction --- at minimum
  \cref{con:greedy-discard} and \cref{con:optimal-discard} --- and the
  two can select different discard sets for the same budget and
  priorities in general. \cref{lem:exhaustion} resolves this specifically
  for \cref{thm:composition-failure}'s own proof, by exhibiting a witness
  relay whose key discard is forced under any construction satisfying
  \cref{def:budget-discard}, but only for stages where the required
  distinctions' cost closely exhausts the local budget
  (\textbf{non-exhausting relays}, below). What remains open is whether
  every \emph{other} claim in this monograph stated to hold for ``every''
  budget relay enjoys the same construction-independence, or whether some
  depend, unexamined, on a specific choice of local discard rule.

  \item \textbf{Non-exhausting relays} (Appendix~\ref{app:proof}).
  \cref{lem:exhaustion} guarantees construction-independent discard only
  for stages where the required distinctions' cost closely exhausts the
  local budget. Whether \cref{thm:composition-failure}'s conclusion is
  robust to choice of construction for relays \emph{without} this
  near-exhaustion property --- where genuine slack remains after funding
  every required distinction --- is open; \cref{con:greedy-discard} and
  \cref{con:optimal-discard} may disagree on which non-required
  distinctions such slack goes to, with unexamined consequences for which
  future tasks such a relay turns out to support.

  \item \textbf{Root-preservation bootstrapping} (Ch.~\ref{ch:budget-relay},
  \cref{rem:root-bootstrap}). Preserving a root object
  (\cref{def:root-object}) so that \cref{prop:root-recovery} remains
  available is itself a repair decision that must be funded before the
  future task space justifying it is known. Whether this can be given a
  more precise decision-theoretic treatment --- for instance, as an
  expected-value calculation over a distribution of possible future tasks
  --- rather than being left at the level of a qualitative design
  recommendation, is open.

  \item \textbf{Infinite relays} (Ch.~\ref{ch:budget-relay},
  \cref{def:substrate-chain}). This monograph treats only substrate
  chains with finitely many stages. Whether an infinite chain behaves
  differently in kind, rather than merely in degree, from the finite
  case --- for instance, whether \cref{prop:monotonicity} or
  \cref{thm:composition-failure} continue to hold, require additional
  hypotheses, or fail outright in the infinite setting --- is
  unaddressed.

  \item \textbf{Branching relays} (Appendix~\ref{app:proof},
  Ch.~\ref{ch:fanout}). Appendix~\ref{app:proof}'s construction is linear
  (\cref{def:substrate-chain}) throughout. Whether an analogous
  construction-independent witness exists for the branching relays of
  Chapter~\ref{ch:fanout}, given the additional complication of
  \cref{prop:silent-disagreement}, has not been checked.

  \item \textbf{Quantitative tightness of \cref{cor:no-free-lunch}}
  (Ch.~\ref{ch:composition-failure}, Appendix~\ref{app:proof}).
  Appendix~\ref{app:proof} establishes only the existence half of this
  corollary --- that cumulative discard is non-decreasing along the
  witnessing construction, realized with equality across arbitrarily many
  pass-through stages --- and explicitly does not establish the stronger
  claim that cumulative discard \emph{risk} grows across relays in
  general, nor whether a sharper bound holds under specific coordination
  mechanisms.

  \item \textbf{General coordination protocols} (Ch.~\ref{ch:coordination},
  Ch.~\ref{ch:synthesis}). Whether a coordination protocol could be
  designed to bound, though not eliminate,
  \cref{thm:composition-failure}'s failure mode across an entire class of
  relays, rather than only via the domain-specific mechanisms of
  Chapter~\ref{ch:coordination}. \cref{prop:pooling-restores}
  (Appendix~\ref{app:comparison}) shows that one particular, more radical
  move --- pooling stage budgets rather than coordinating separate ones,
  per Section~\ref{sec:coord-vs-pooling}'s distinction --- can restore
  sufficiency in a specific witness; it does not answer this broader
  question, since pooling eliminates the stage-separation the question is
  asking about coordinating \emph{within}, rather than supplying a
  coordination protocol in the sense of Chapter~\ref{ch:coordination}'s
  taxonomy.

  \item \textbf{The closing synthesis} (\cref{conj:closing},
  Ch.~\ref{ch:synthesis}). Labeled a conjecture because the unification of
  \cref{thm:composition-failure} with the fidelity continuum has not
  itself been given an independent proof beyond what its components
  separately establish.

  \item \textbf{Fidelity and Transport Stability}
  (\cref{rem:fidelity-transport-open}, Ch.~\ref{ch:four-operations};
  Ch.~\ref{ch:fidelity-continuum}). Whether $\fidelity{r}{\tau}$
  (\cref{def:fidelity}) can be derived from, bounded by, or otherwise shown
  compatible with Distinguishability Geometry's Transport Stability theorem
  (\cref{thm:transport-stability}) is unexamined; the two are currently
  related only by analogy.

  \item \textbf{Revision under a changing state space --- resolved}
  (\cref{prop:relay-replacement-paradigm-shift}, Ch.~\ref{ch:four-operations};
  Ch.~\ref{ch:science}). Previously open: whether relay revision under a
  changing state space could be expressed formally at all, given that
  Distinguishability Geometry's revision operation requires a fixed state
  space. Resolved by Observability Is Restorability's Paradigm Shift
  Theorem (\cref{thm:paradigm-shift}), which characterizes exactly this
  case as repair-repertoire discontinuity producing incomparable
  restorable-distinction sets. What remains open is only the
  interpretive question of which of science's historical episodes fall
  under the fixed-state-space case
  (\cref{prop:relay-replacement-partial-revision}) versus this one, which
  this monograph does not attempt to adjudicate.

  \item \textbf{Quantitative deficit bounds for this monograph's discard sets}
  (\cref{prop:rendering-is-projection}, Ch.~\ref{ch:four-operations}).
  Every discard set in this monograph (\cref{def:discard-set}) could in
  principle carry a quantitative lower bound via Distinguishability
  Geometry's Projection Deficit Bound (\cref{thm:projection-deficit-bound}),
  $\Delta\delta(r) \ge H(X \mid Y)$. No chapter has computed this bound for
  any of its worked instances; the discard-set apparatus throughout this
  book remains qualitative (present/absent) rather than quantitative.

  \item \textbf{The measure of a persistence set}
  (\cref{rem:root-recovery-persistence-set}, Ch.~\ref{ch:four-operations};
  Ch.~\ref{ch:budget-relay}). Root recovery (\cref{prop:root-recovery})
  establishes that a specific transformation lies in a persistence set
  $P(d)$ (\cref{def:persistence-set}) when a root object is preserved. What
  $\mu(P(d))$ is in general --- what fraction of the relevant transformation
  family admits recovery, and at what cost it can be enlarged --- has not
  been computed, and would give the root-preservation bootstrapping problem
  above a quantitative rather than qualitative target.
  \cref{rem:witnesses-persistence-measure} records a specific candidate
  from a weaker-tier source, The Geometry of Witnesses's reconstruction
  redundancy $\rho_W(r) = |C(r)|$, which is unverified but concrete enough
  to be worth checking directly rather than left as a placeholder for
  ``some measure or other.''

  \item \textbf{One remaining unverified import, one retraction}
  (Front matter, Relation to Prior Work). \emph{The Geometry of
  Reconstruction} was cited in an earlier draft and has been retracted: no
  such text has been located, and the idea it was cited for (second-order
  repair operators maintained by witnesses) is now attributed instead to
  \emph{The Geometry of Witnesses}, a text the author reports as
  forthcoming or separately authored. This replacement citation has
  itself not been independently read or confirmed as existing in completed
  form, and should be treated as no more reliable than the citation it
  replaced until it is checked directly.

  \item \textbf{The transcription tier} (Front matter, Relation to Prior
  Work; \S\S\ref{sec:abt}--\ref{sec:goc}). \emph{Admissibility Before
  Truth} and \emph{The Geometry of Correction} are integrated on the
  strength of formulas and theorem statements the author transcribed
  directly from source PDFs, marked \transcribed{} throughout rather than
  \imported{}. This is stronger evidence than the paraphrase-level
  citations these two sources previously carried, but it is not the same
  as an independent reading of the primary text, which remains the
  standard met only by \emph{Conservation of Ambiguity}, Distinguishability
  Geometry, Persistence Before Truth, and Observability Is Restorability.
  Whether the transcribed formulas are complete and correctly contextualized
  relative to their source documents has not been independently checked.

  \item \textbf{The admissibility collision, sharpened but not resolved}
  (\S\ref{sec:abt}, Ch.~\ref{ch:four-operations}). Four distinct formal
  notions bearing the name ``admissibility'' are now on record: this
  monograph's epistemic $\admis{\tau}{r}$, the activation-manifolds
  application paper's $A: M \to \mathbb{R}_{\ge 0}$, Observability Is
  Restorability's observer-relative $\mathcal{A}_O$, and Admissibility
  Before Truth's domain restriction $\mathrm{dom}(\mathcal{T}) \subseteq
  \mathrm{Rest}(\mathcal{R})$. \cref{prop:sufficiency-presupposes-evaluability}
  connects the last of these directly to this monograph's own apparatus by
  a stated proof, which is more traction than the other two comparisons
  have yielded, but reconciling all four, or the primary Admissibility
  Field text none of them may fully represent, remains undone.

  \item \textbf{Repair-interference formula discrepancy, internal to the
  corpus} (\cref{rem:interference-inconsistency}, Ch.~\ref{ch:four-operations}).
  \emph{Conservation of Ambiguity} attributes its derivative-form repair
  interference functional to Observability Is Restorability; Observability
  Is Restorability's own verified text defines interference algebraically
  instead ($I(d_i,d_j) = C_R(d_i \wedge d_j) - C_R(d_i) - C_R(d_j)$). This
  monograph uses the derivative form consistently and correctly attributes
  it to \emph{Conservation of Ambiguity}, but the discrepancy between the
  two source texts themselves is unresolved and not this monograph's to
  settle.

  \item \textbf{When reconciliation resolves automatically versus requires
  adjudication} (\cref{rem:reconciliation-conflict}, Ch.~\ref{ch:fanout}).
  Reconciliation (\cref{def:reconciliation}) recovers what non-overlapping
  branch discard destroyed, but a genuine conflict -- both branches
  preserving the same distinction with incompatible values -- is not
  resolved by the intersection rule and falls back on
  Chapter~\ref{ch:coordination}'s mechanisms instead. Whether the boundary
  between these two cases can be characterized precisely in advance,
  rather than discovered only by attempting a merge and finding a
  conflict, is open.

  \item \textbf{Coordinating the coordinators} (\S\ref{sec:witnesses},
  Ch.~\ref{ch:coordination}). Chapter~\ref{ch:coordination}'s taxonomy
  coordinates distinctions across separated stages but is silent about
  what maintains the coordinating mechanism itself -- the shared ledger,
  the standing court, the augmentation policy. This is precisely the
  question \S\ref{sec:witnesses} reports The Geometry of Witnesses as
  raising one level up, at that source's weaker, unverified tier; this
  monograph does not develop a coordinated-coordinator layer of its own,
  and the gap between coordinating distinctions and coordinating the
  capacity to coordinate them remains open.
\end{itemize}
